\documentclass{report}
\usepackage{amsmath,amssymb}
\setlength{\parindent}{0mm}
\setlength{\parskip}{1em}
\begin{document}
\begin{center}
\rule{6in}{1pt} \
{\large Ming Yan \\
{\bf Binary Matching Pursuit for Impulse Noise Removal}}

520 Portola Plaza \\ LOS ANGELES \\ California \\ United States \\ 90095
\\
{\tt yanm@ucla.edu}\end{center}

This article studies the problem of image restoration of observed images
corrupted by impulse noise and other types of noise (e.g. zero-mean
Gaussian white noise). Since the pixels damaged by impulse noise contain
no information about the true image, the damaged pixels can also be
considered as missing information. If the pixels corrupted by impulse
noise are known, then the image restoration becomes a standard image
inpainting problem. However, the set of damaged pixels is usually
unknown, thus how to find this set correctly is a very important problem.
We proposed a method using binary matching pursuit that can
simultaneously find the damaged pixels and restore the image. This method
can also be applied to situations where the damaged pixels are not
randomly chosen, but follow some unknown procedure. By iteratively
restoring the image and updating the set of damaged pixels, this method
has better performance than other methods, as shown in the experiments.


\end{document}
